\section{Implementation notes}
\label{app:impl}

\subsection{Training failures that shaped the implementation}
\label{app:failures}

Three failures were reproducible and are worth recording, because each one
produces a network that looks superficially fine but cannot be certified. The
runs behind them are kept in \url{runs/appendix_fourier2d/} and
\url{runs/appendix_failures/}; the Fourier and hard-mask runs are outputs of
\texttt{scripts/exp6\_certificate\_training\_2d.py}, the soft-boundary and
unnormalized runs of \url{scripts/appendix_failures.py}.

\paragraph{Random Fourier features in two dimensions} The one-dimensional
runs use random Fourier features \citep{tancik2020} and profit from them. In
two dimensions the same configuration (scale $2.5$) stalls: at $6000$ epochs,
on the harmonic potential whose discrete eigenvalue is $1.0004$, it reaches
$E=1.5624$ and overlap $0.877$ (the $1500$-epoch snapshot is $E=1.5636$ and
$0.877$) against $E=1.0004$ and overlap $0.99997$ for the plain MLP, and the
single-state certificate refuses the output with a residual of $1.17$
(\texttt{runs/appendix\_fourier2d/}). Each feature
$\sin/\cos(\omega_1x+\omega_2y)$ must cover a direction in the plane, and the
sample efficiency of a fixed random direction set degrades with dimension. The
2D driver therefore disables Fourier features by default. This is an empirical
configuration choice, not a claim about Fourier features in general.

\paragraph{The zero function in an unnormalized loss} If the loss is the
residual of $\psi$ itself rather than of the unit vector $\psi/\norm{\psi}$,
the zero function is a global minimum and not merely a stationary point: the
Rayleigh quotient is scale invariant, so the residual is exactly linear in
$\psi$ and $L(c\psi)=c^2L(\psi)$, whence $L(0)=0\le L(\psi)$. No term then
fixes the scale: the amplitude stays wherever the initialization put it, and
any tolerance can be met by shrinking rather than by improving the shape. In a
$3000$-epoch run at $N=41$ the field ended at
$\norm{\psi}^2=6.0\times10^{-5}$, a maximum amplitude of $1.5\times10^{-2}$ (rms
$7.5\times10^{-3}$ over the grid), with
a shape that is already accurate (overlap $1.0000$ after normalization, energy
error $1.1\times10^{-5}$). Measured on the field as returned, without
normalization, the residual norm is $2.9\times10^{-4}$ while the $L^2$
distance from the discrete ground state is $0.99$---a residual three thousand
times smaller than the error it is supposed to measure
(\texttt{scripts/residual\_vs\_error.py}). Every diagnostic
computed on the normalized field therefore looks excellent while the field
itself carries a scale that the loss does not determine, and the training loss
then measures the amplitude rather than the error.
Fixing the scale by an explicit term---the penalty $(\norm{\psi}_h^2-1)^2$ in
one dimension, division by the norm in two (equation~\eqref{eq:loss})---removes
the scale direction and excludes the zero solution; the denominator is clamped
at $10^{-30}$ to keep the gradient finite during the first steps.

\paragraph{Soft boundary penalties in two dimensions} A boundary penalty with
weight $100$ does not enforce a Dirichlet condition to anything like machine
precision, because the constrained nodes are an $O(1/N)$ fraction of the grid
while the five-point Laplacian amplifies their error by $h^{-2}$. After
$3000$ epochs at $N=41$ on the box potential the penalty leaves the residual
at $5.72$ and the energy error at $5.0\times10^{-2}$ ($E=9.9144$ against the
discrete $9.8645$), with boundary nodes carrying values up to $0.22$. The
certificate still returns \texttt{valid}---its hypotheses hold for whatever
field it is given---but the enclosure is $4.68$ wide against a true error of
$5.0\times10^{-2}$, more than five orders of magnitude looser than on a
hard-constrained field. Multiplying the network output by the mask
$M=(x-a)(b-x)(y-a)(y-b)$ enforces the boundary condition exactly and is a
statement of the problem, not of the solution; with everything else unchanged
the same configuration reaches residual $1.2\times10^{-2}$, energy error
$1.2\times10^{-6}$, enclosure width $2.0\times10^{-5}$ and angle bound
$9.0\times10^{-4}$ (\texttt{runs/appendix\_failures/}). The mask costs
expressivity near the boundary, which is irrelevant for an eigenfunction that
vanishes there.

\subsection{Interval construction}
\label{app:algorithm}

\begin{center}
\fbox{\parbox{0.93\linewidth}{\small
\textbf{Algorithm 1} Verified or refused certificate for $(\psi,\text{grid},k^\ast)$.
\begin{enumerate}
\item[1.] $\eta\leftarrow\eta(\psi)$, $r\leftarrow\Hh\psi-\eta\psi$,
  $\rho\leftarrow\norm{r}_h$. If $\rho=0$ return the exact certificate.
\item[2.] $k\leftarrow N(\eta)$ (verified count). Obtain an approximate
  spectrum from ARPACK for states $k-1,k,k+1$ only.
\item[3.] For \texttt{state} in $\{k-1,k\}$: build
  $a=\hat\lambda_{\text{state}-1}+0.1(\hat\lambda_{\text{state}}-\hat\lambda_{\text{state}-1})$
  and $b=\hat\lambda_{\text{state}+1}-0.1(\hat\lambda_{\text{state}+1}-\hat\lambda_{\text{state}})$;
  require $a<\eta<b$ and $N(b)-N(a)=1$; require
  $\min(\eta-a,b-\eta)>10^{-9}s$; keep the candidate with the larger gap.
\item[4.] If no candidate verifies, run the count-only bisection fallback and
  verify its interval in the same way; if that fails, return
  \texttt{valid=False} with the reason.
\item[5.] Compute $\gamma=\min(\eta-a,b-\eta)$. If
  $\rho^2<(b-\eta)(\eta-a)$ return the Temple enclosure \eqref{eq:temple};
  otherwise return the Weinstein interval and mark the Temple bound
  unavailable. In both cases return $\arcsin(\min(1,\rho/\gamma))$.
\end{enumerate}}}
\end{center}

\subsection{Reproducibility details}
\label{app:repro}

All runs were performed on a laptop CPU (Windows 11, Python 3.11,
NumPy~2.4 \citep{harris2020}, SciPy~1.17 \citep{virtanen2020}, PyTorch~2.12,
the thread count fixed per run in the sweep manifest---four for a
one-dimensional run, two for a two-dimensional one). The JSON manifests record
the wall-clock time of each job. The reference spectra used for validation are
computed by eigenvalue solvers independent of the certificate, applied to the
same assembled operator: in one dimension
by LAPACK's symmetric tridiagonal solver \citep{virtanen2020}, with the
eigenvalue taken as the Rayleigh quotient of the returned vector so that the
absolute accuracy is not limited by $\varepsilon\norm{\Hh}$ ($7\times10^{-12}$
to $2\times10^{-9}$ on the three one-dimensional grids); in
two dimensions by a dense LAPACK decomposition of the full
$6241\times6241$ matrix (\texttt{scripts/dense\_reference\_2d.py}), which
agrees with the ARPACK values to all ten printed digits. ARPACK is used only
to propose interval endpoints, never to accept them.

The test suite contains 143 tests and runs with \texttt{pytest} or with an
included minimal runner when \texttt{pytest} is unavailable. The tests that
matter for the correctness of the certificate are the ones that compare the
counts with dense diagonalization (both dimensions), the ones that assert
containment for exact eigenvectors, the ones that assert \emph{refusal} of a
degenerate state by the single-state certificate, and
(\texttt{test\_certify\_subspace.py}) the ones that certify the same
degenerate pair as a two-dimensional subspace, check that the certified angle
bound dominates the true principal angle, and check that the certificate is
invariant under a change of basis of the trial subspace.
